\section{Communication over persistent financial stakes}
\label{sec:conduct}

\subsection{Bargaining limits and changing cash needs}
Cash affects both the willingness to transact and what a transaction enables. Among proposer-only completion requests, holders below \$100 reject 17/69; holders at \$300 or more reject 120/175. Most low-cash responses are counters, 44/69. The pooled difference includes later negotiations and Opus's response style. Of the 69 low-cash offers, 43 are later offers and 21 face Opus. Gemini remains resistant even at low cash, rejecting ten of eleven such requests. There is no uniform threshold at which a holder releases a rival's completing deed.

Longer histories show positions and terms changing together. Fourteen completing transfers follow an earlier refusal or counter by the same holder to the same prospective completer. Pairing settlement with the first earlier resistance on a previous turn, cash falls in thirteen, from a median \$571.50 to \$82. But the net cash offered rises in eight and only five pairs preserve the same reciprocal deed bundle. These histories document the course of negotiation, not an effect attributable solely to liquidity.

\TranscriptTitle{g7}{A settled bluff and a later financing bargain}
Sol asks \$625 for North Carolina. Grok replies:
\TranscriptQuote{g7_d42_public}
Its private note retains the higher contingency:
\TranscriptQuote{g7_d42_private}
Sol accepts \$600. The public limit understates an explicit private willingness to pay, making this a documented bluff under the task's permitted conduct. It contrasts with Opus's refused final-price claim in G2.

Later in G7, Sol buys Park Place, builds three dark-blue houses and has \$84 remaining. Gemini's note identifies a possible financing opportunity:
\TranscriptQuote{new-g7-d141-timed-bid}
Gemini bids \$350 for New York; Sol asks \$700 to fund its own blues while limiting the buyer's oranges. Gemini pays, mortgages five deeds and builds four orange houses that turn. Sol adds three dark-blue houses three turns later, then sells one on the next turn to meet rent. Gemini collects \$270 orange rent before its elimination by Sol. The bargain finances competing investments and subsequent liquidity demands. Earlier offers had different property terms, so the sequence does not isolate Sol's cash as the reason agreement occurred.
\TranscriptSource{G7 d41--43 and d141 onward. Public exchange and subsequent transactions, mortgages and builds are traced in the saved records.}

\subsection{Funding a rival to preserve the competitive position}
Ownership makes a rival's survival relevant to other players. A player may prefer leaving a group with a weak holder to seeing it inherited by a stronger creditor. Economic theories of sales with identity-dependent externalities formalize why the identity of an owner can matter to nonbuyers~\cite{jehiel1996externalities}. The following incident supplies explicit statements about that consequence.

\TranscriptTitle{g13}{Selective presentation in an inheritance pitch}
Grok holds \$35 and stands on Atlantic before Sol's Short Line. It offers mortgaged Pennsylvania to Opus for \$200:
\TranscriptQuote{g13-opening-insurance}
Opus counters at \$130. Grok's next applied note recognizes that this payment suffices if it also mortgages Ventnor, while its provider-returned explanation says:
\TranscriptQuote{g13-omission-explanation}
Opus holds \$1,926, accepts \$200 and pays \$16 transfer interest. Its explanation considers two possible recipients of a future estate:
\TranscriptQuote{integrated-g13-buyer}
The buyer weighs both denying Sol the yellows and potentially inheriting them itself. Its solvency reasoning already appears in the \$130 counter, after Grok's opening appeal and before the omission-framed reply. The records support selective presentation of a financing option and uptake of the estate framing. They do not show that the omitted mortgage caused the extra \$70 payment or that Opus believed the claimed minimum.

Grok misses Short Line on the next roll and later goes bankrupt there at turn 150, transferring yellow to Sol. Opus completes green at auction seven turns after the purchase and ultimately wins. The pitch's 5/36 figure repeats the distance miscount in Grok's contemporaneous note; the later note uses the correct 4/36 for a direct roll of nine. Neither number describes all movement within a turn with doubles. The recorded omission concerns collateral, while the pitch's probability error has no comparable evidence of deliberate falsification.
\TranscriptSource{G13 d250--253 and d311. Opus uses recovery. The omission is in Grok's provider-returned explanation; the public pitch is present in Opus's request.}

\subsection{Financial denial and contingent assurances}
\TranscriptTitle{g20}{A coordination appeal with unilateral follow-through}
Opus sells Illinois to Grok for \$500, completing reds and leaving Grok \$222. Gemini addresses the other players:
\TranscriptQuote{g20-coordination-appeal}
Sol had privately declined to buy Grok's New York one turn before the appeal because the cash would finance dangerous development. Grok's next note recognizes Gemini's appeal. It collects another \$235, builds one house on each red, leaving \$7, and mortgages New York for \$100 before offering Vermont to Gemini for \$150:
\TranscriptQuote{g20-concrete-refusal}
Grok mortgages Vermont for \$50 instead. Opus subsequently offers cash for Grok's deeds, explicitly as finance for more red houses, but no sale settles. Grok is eliminated first and Gemini wins. The incident contains a public attempt to coordinate and the speaker's corresponding refusal. Another player's contrary offer prevents treating it as established collective compliance or a functioning cartel.
\TranscriptSource{G20 d95--102, d129--130 and d156--160. The appeal is in Grok's d97 prompt.}

\TranscriptTitle{g18}{Withdrawing an assurance to the last opponent}
Only Opus and Sol remain. Opus bids \$200 and then \$250 for Sol's mortgaged Pacific, which would complete its greens, and Sol refuses. Opus publicly restates the standing offer four times. The fourth follows a \$750 rent:
\TranscriptQuote{g18-assurance}
The explanation returned with that same renewal already records a reservation:
\TranscriptQuote{integrated-g18-reservation}
Another \$750 rent forces Sol to liquidate and leaves \$49. After larger packages fail, Sol invokes the standing offer:
\TranscriptQuote{g18-invocation}
\TranscriptQuote{g18-withdrawal}
Opus's returned explanation objects to prolonging a nearly bankrupt rival and to the cost of green development. It does not explicitly invoke inheritance. Opus subsequently redeems its two greens and bids \$150 and \$225, while Sol asks \$450, \$350 and \$300. Nothing settles. Opus eventually wins at Sol's elimination, whose terminal estate has no later earning opportunity.

The public renewal and returned reservation coexist before Sol's further financial deterioration. The subsequent withdrawal makes that reservation consequential for the negotiation. The first \$250 bid reads as a live offer in its returned explanation; the record does not show whether the renewed assurance would have been honored under other conditions or a cost incurred by Sol through reliance.
\TranscriptSource{G18 d421--461. Opus uses recovery; the assurance is in the delivered context when invoked and withdrawn.}

Taken together, the incidents identify specific communicative acts: G7 juxtaposes a public limit and a private contingency; G13 explicitly describes an omission; G19 corrects a false premise while accepting the terms; G20 seeks collective financial denial; G18 renews an assurance while expressing a reservation, then withdraws it.
